\documentclass[11pt]{article}
\usepackage[T1]{fontenc}
\usepackage{lmodern}
\usepackage[margin=1in]{geometry}
\usepackage{amsmath,amssymb,amsthm,mathtools}
\usepackage{booktabs,array,tabularx}
\usepackage{float}
\usepackage{graphicx}
\usepackage{microtype}
\usepackage[hidelinks]{hyperref}
\usepackage{xcolor}

\newtheorem{theorem}{Theorem}
\newtheorem{lemma}[theorem]{Lemma}
\newtheorem{corollary}[theorem]{Corollary}
\newtheorem{proposition}[theorem]{Proposition}
\newtheorem{remark}[theorem]{Remark}
\usepackage[nameinlink,noabbrev]{cleveref}
\crefname{theorem}{Theorem}{Theorems}
\Crefname{theorem}{Theorem}{Theorems}
\crefname{lemma}{Lemma}{Lemmas}
\Crefname{lemma}{Lemma}{Lemmas}
\crefname{corollary}{Corollary}{Corollaries}
\Crefname{corollary}{Corollary}{Corollaries}
\crefname{proposition}{Proposition}{Propositions}
\Crefname{proposition}{Proposition}{Propositions}
\crefname{remark}{Remark}{Remarks}
\Crefname{remark}{Remark}{Remarks}
\newcommand{\R}{\mathbb{R}}
\newcommand{\Vol}{\operatorname{Vol}}
\newcommand{\Lip}{\operatorname{Lip}}

\title{A Certified Lower Bound for the Maximum Volume\\Wrappable by a Disk}
\author{Sho Hoshino}
\date{August 11, 2026}

\begin{document}
\maketitle

\begin{abstract}
For the closed disk $D_R\subset\R^2$, let
\[
W(R):=\sup\{\Vol(B): B\text{ is wrappable by }D_R\}
\]
in the $1$-Lipschitz wrapping formalization of Nandakumar.  We prove the explicit certified lower bound
\[
W(R)\ge C R^3,
\qquad
C=0.17984349416291615\ldots,
\]
where $C$ is an exact rational number.  More precisely, we give a $17$-vertex, $30$-face embedded nonconvex polyhedral sphere and a continuous $1$-Lipschitz map from $D_R$ onto it.  Two independent exact-arithmetic programs certify strict nonexpansion on all $30$ source triangles, the topology of the quotient complex, all $435$ triangle-pair intersections, nonconvexity, and the exact enclosed volume.

Nandakumar's Question~1 asks only whether $D_\pi$ can wrap some nonconvex solid of volume greater than $4\pi/3$.  Under the stated $1$-Lipschitz formalization, that question has a short soft affirmative answer: an oblate spheroid with equatorial radius $23/20$ and polar radius $4/5$ has volume greater than the unit ball and north--south meridian length strictly less than $\pi$; a sufficiently small inward rotational notch preserves both strict inequalities while making the body nonconvex.  A simple radial lemma then gives a $1$-Lipschitz wrap by $D_\pi$.  Thus the point of the certified polyhedral construction is not existence of a nonconvex improvement over the sphere, but the substantially stronger quantitative lower bound above.  For comparison, the classical Mylar balloon is the sharp convex body-of-revolution benchmark with coefficient $0.152315527\ldots$, so our certified coefficient is about $18.07\%$ larger.
\end{abstract}

\section{Introduction}
Nandakumar recently formulated a volume-maximization problem for wrapping three-dimensional solids by a planar sheet that may be folded or wrinkled but not stretched or torn~\cite{Nandakumar2026}.  In his formalization, a planar region $D$ wraps a compact connected solid when its boundary surface is realized by a $1$-Lipschitz image of $D$, with the boundary of $D$ glued to itself in the image.  For the disk $D_R$ of radius $R$, it is therefore natural to write
\begin{equation}\label{eq:Wdef}
W(R):=\sup\{\Vol(B):B\text{ is wrappable by }D_R\}.
\end{equation}
Scaling gives $W(R)=R^3W(1)$ whenever the supremum is interpreted in this fixed formalization, so the main quantitative problem is to understand the dimensionless constant $W(1)$.

Nandakumar's main conjecture is explicitly nonconvex: every wrappable convex solid should be beaten in volume by some wrappable nonconvex solid~\cite[Claim~1]{Nandakumar2026}.  His Question~1 fixes the unit ball as a test case and asks whether $D_\pi$ can wrap a nonconvex solid of volume greater than $4\pi/3$~\cite[Question~1]{Nandakumar2026}.  A subsequent MathOverflow remark by Nandakumar already records convex wrappable bodies with volume larger than the unit sphere~\cite{NandakumarMO2026}.  Under the stated $1$-Lipschitz formalization, one can go one step further: as shown in \cref{sec:soft}, a tiny inward dent of a suitably chosen oblate spheroid gives a short elementary affirmative answer to Question~1.  Thus Question~1 itself is not the quantitative difficulty addressed here.

The purpose of this note is to give a substantially stronger explicit lower bound for \eqref{eq:Wdef}.  The construction is finite and completely specified: we give the source triangulation, target coordinates, and piecewise-affine map.  The proof then reduces to a finite exact certificate.  Two supplied implementations use different arithmetic and different triangle-intersection algorithms.

Our main result is the following.

\begin{theorem}[Certified disk-wrapping lower bound]\label{thm:main}
Let $D_R\subset\R^2$ be the closed disk of radius $R>0$.  There exists a continuous $1$-Lipschitz map
\[
F_R:D_R\longrightarrow\R^3
\]
whose image is an embedded nonconvex polyhedral $2$-sphere $\partial B_R$.  The bounded solid $B_R$ has volume
\begin{equation}\label{eq:Cexact}
\Vol(B_R)=C R^3,
\end{equation}
where
\[
C=
\frac{
359686988325832312742347877865122384142690317017149
}{
2000000000000000000000000000000000000000000000000000
}
=0.17984349416291615\ldots .
\]
In particular,
\[
W(R)\ge C R^3>0.179843R^3.
\]
\end{theorem}

The classical Mylar balloon provides a useful convex benchmark.  Among convex bodies of revolution with intrinsic diameter $R$, it has the sharp volume
\[
V_{\rm Mylar}(R)=0.152315527\ldots R^3,
\]
and it is itself wrappable by the single disk $D_R$.  Hence \cref{thm:main} improves this axisymmetric convex benchmark by about $18.07\%$.  We do not claim that Mylar is optimal among all convex bodies wrappable by $D_R$, nor do we claim to determine $W(R)$.

The geometry of the certified witness was discovered by numerical search, but the theorem depends only on the final rational-coordinate certificate.  Its boundary identification has one short pleat branch and otherwise consists of doubled seam segments; no physical analogy is needed in the proof.

\section{A soft affirmative answer to Nandakumar's Question~1}\label{sec:soft}
We first isolate a simple observation that applies to convex and nonconvex surfaces of revolution alike.

\begin{lemma}[Radial wrapping of a surface of revolution]\label{lem:radial}
Let
\[
\gamma(s)=(\varrho(s),z(s)),\qquad 0\le s\le L,
\]
be a simple rectifiable arc in the half-plane $\varrho\ge0$, parametrized by arclength, with
\[
\varrho(0)=\varrho(L)=0,
\qquad
\varrho(s)>0\quad(0<s<L).
\]
Let $S$ be the embedded surface obtained by revolving $\gamma$ about the $z$-axis.  If $L\le R$, then $D_R$ wraps the bounded solid enclosed by $S$ in Nandakumar's $1$-Lipschitz sense.
\end{lemma}

\begin{proof}
Set $q=L/R\le1$.  In polar coordinates $(r,\theta)$ on $D_R$, define
\[
F(r,\theta)=
\bigl(\varrho(qr)\cos\theta,\,\varrho(qr)\sin\theta,\,z(qr)\bigr).
\]
Since $\gamma$ is parametrized by arclength,
\[
\varrho'(s)^2+z'(s)^2=1
\]
for almost every $s$, so $|\varrho'|\le1$.  Because $\varrho(0)=0$,
\[
0\le\varrho(s)\le s.
\]
Away from the rotation axis the pullback metric is
\[
F^*g=q^2\,dr^2+\varrho(qr)^2\,d\theta^2
\le q^2\bigl(dr^2+r^2d\theta^2\bigr)
\le dr^2+r^2d\theta^2.
\]
Integrating lengths of curves and using continuity at the axis gives $\Lip(F)\le q\le1$.  The map is onto $S$, while the whole boundary circle $r=R$ is sent to the endpoint $\gamma(L)$ on the axis.  Hence the disk boundary is self-glued and the image is the closed surface $S$.
\end{proof}

\begin{proposition}[A soft answer to Question~1]\label{prop:softQ1}
A disk of radius $\pi$ wraps a nonconvex solid of volume strictly greater than $4\pi/3$.
\end{proposition}

\begin{proof}
Consider the oblate spheroid with equatorial radius
\[
a=\frac{23}{20}
\]
and polar radius
\[
c=\frac45.
\]
Its volume is
\[
V_0=\frac{4\pi}{3}a^2c
=\frac{4\pi}{3}\frac{529}{500}
>\frac{4\pi}{3}.
\]
A north--south generating meridian is
\[
\gamma_0(t)=(a\sin t,c\cos t),\qquad 0\le t\le\pi.
\]
If $L_0$ denotes its length, Cauchy--Schwarz gives
\begin{align*}
L_0^2
&=\left(\int_0^\pi
\sqrt{a^2\cos^2t+c^2\sin^2t}\,dt\right)^2\\
&\le \pi\int_0^\pi
\bigl(a^2\cos^2t+c^2\sin^2t\bigr)\,dt\\
&=\pi^2\frac{a^2+c^2}{2}
=\pi^2\frac{157}{160}<\pi^2.
\end{align*}
Thus $L_0<\pi$ with strict slack.

We now make an arbitrarily small inward notch near the equator.  For $\delta>0$ small, put
\[
h=c\sin\delta,
\qquad
p_\pm=(a\cos\delta,\pm h).
\]
Replace the elliptical arc between $p_+$ and $p_-$ by the two line segments joining these points to
\[
v=(a\cos\delta-\varepsilon,0),
\qquad
0<\varepsilon<a\cos\delta.
\]
Because the ellipse is convex, the two new segments lie inside it.  Revolving the modified simple meridian gives an embedded solid $B_{\delta,\varepsilon}$ obtained from the spheroid by a rotational inward notch.

Its meridian length satisfies
\[
L_{\delta,\varepsilon}
\le L_0+2\sqrt{h^2+\varepsilon^2},
\]
so for sufficiently small $\delta$ and $\varepsilon$ we still have $L_{\delta,\varepsilon}<\pi$.  The removed volume lies inside the slab $|z|\le h$ and the cylinder $\varrho\le a$, hence
\[
0\le V_0-\Vol(B_{\delta,\varepsilon})
\le 2\pi a^2h.
\]
Since $h\to0$ with $\delta\to0$ and $V_0>4\pi/3$, the volume remains greater than $4\pi/3$ for sufficiently small $\delta$.

Finally, the meridional cross-section is nonconvex: it contains $p_+$ and $p_-$, but their midpoint $(a\cos\delta,0)$ lies outside the notched body because the new boundary radius at $z=0$ is $a\cos\delta-\varepsilon$.  Therefore $B_{\delta,\varepsilon}$ is nonconvex.  Applying \cref{lem:radial} with $R=\pi$ completes the proof.
\end{proof}

Thus the qualitative content of Question~1 follows without computer assistance.  The rest of the paper concerns the stronger quantitative problem of pushing the coefficient in \eqref{eq:Wdef} upward with a fully explicit certificate.

\section{A convex benchmark}\label{sec:context}
The classical Mylar balloon is the convex surface of revolution arising from the fixed-meridian-length volume-maximization problem studied by Paulsen and subsequently by Mladenov and Oprea~\cite{Paulsen1994,MladenovOprea2003}.  For a convex body of revolution, the intrinsic diameter equals the north--south meridian length $R$.  Indeed, in meridional arclength coordinates $(s,\theta)$ the surface metric has the form $ds^2+\varrho(s)^2d\theta^2$, with $0\le s\le R$.  Any two points at arclengths $s_1,s_2$ can be joined through one of the poles with length at most
\[
\min\{s_1+s_2,\,2R-s_1-s_2\}\le R,
\]
whereas every path from the north pole to the south pole has length at least $\int|ds|\ge R$.  Thus fixing intrinsic diameter is equivalent, in this convex axisymmetric class, to fixing the generating meridian length.

For completeness, the sharpness of the Mylar profile in this class admits a short global argument.  Parametrize the meridian by arclength and write it as $(\varrho(s),z(s))$.  For a general convex meridian these coordinate functions are Lipschitz; throughout this paragraph derivatives are understood almost everywhere.  Convexity makes $z$ monotone and $\varrho$ increase from $0$ to its maximum $a$ and then decrease to $0$.  Hence
\[
\frac{V}{\pi}=\int_0^R \varrho^2\sqrt{1-\varrho'^2}\,ds.
\]
For every $0\le\varrho\le a$, Cauchy--Schwarz gives
\[
\varrho^2\sqrt{1-\varrho'^2}
+|\varrho'|\sqrt{a^4-\varrho^4}
\le a^2.
\]
After integration and the monotonicity of the two halves of the meridian,
\[
\frac{V}{\pi}
\le a^2R-2\int_0^a\sqrt{a^4-u^4}\,du
=a^2R-2Ka^3,
\qquad
K:=\int_0^1\sqrt{1-u^4}\,du.
\]
The right-hand side is maximized at $a=R/(3K)$.  Since $K=2I/3$ for the constant $I$ below, this is $a=R/(2I)$.  Equality in Cauchy--Schwarz requires, on each monotone half of the meridian,
\[
|\varrho'|=\sqrt{1-(\varrho/a)^4},
\qquad
\sqrt{1-\varrho'^2}=\frac{\varrho^2}{a^2}.
\]
Integrating the first relation from a pole to the equator and back gives
\[
R=2a\int_0^1\frac{du}{\sqrt{1-u^4}}=2aI,
\]
so the equality profile has exactly the maximizing value $a=R/(2I)$.  This is the classical Mylar profile, and hence the upper bound is attained.

Let
\[
I=\int_0^1\frac{du}{\sqrt{1-u^4}}
=\frac{\sqrt\pi\,\Gamma(1/4)}{4\Gamma(3/4)}.
\]
The Mylar volume is
\begin{equation}\label{eq:mylar}
V_{\rm Mylar}(R)
=\frac{\pi}{12I^2}R^3
=\frac43\left(\frac{\Gamma(3/4)}{\Gamma(1/4)}\right)^2R^3
=0.152315527\ldots R^3.
\end{equation}
Its generating meridian has length $R$, so \cref{lem:radial} gives a direct $1$-Lipschitz wrap by the single disk $D_R$.  This uses the Mylar balloon only as a target body; it does not identify the source with the doubled circular sheets used in the classical physical construction.

For comparison, the sphere of radius $R/\pi$ also has north--south meridian length $R$, and hence is wrapped by $D_R$ by the same lemma.  Its volume is
\[
V_{\rm sph}(R)=\frac{4}{3\pi^2}R^3=0.1350949115\ldots R^3.
\]
Thus the explicit benchmarks are
\begin{center}
\begin{tabular}{@{}lc@{}}
\toprule
body & volume divided by $R^3$\\
\midrule
sphere & $0.1350949115\ldots$\\
sharp convex body-of-revolution Mylar benchmark & $0.1523155270\ldots$\\
certified nonconvex body of this paper & $0.1798434942\ldots$\\
\bottomrule
\end{tabular}
\end{center}
The certified construction therefore exceeds the Mylar benchmark by about $18.07\%$.  This is only a benchmark within the convex axisymmetric class; no optimality among all convex wrappable bodies is asserted.

\section{The source complex and the one-pleat seam}\label{sec:source}
It suffices to construct the case $R=1$; scaling source and target by $R$ then multiplies volume by $R^3$.

Let $H$ be the regular decagon inscribed in the unit circle, with outer vertices
\[
o_k=(\cos(2\pi k/10),\sin(2\pi k/10)),\qquad k=0,\ldots,9.
\]
Set
\[
\rho=\frac{13}{20},\qquad i_k=\rho o_k,
\]
and let $z=(0,0)$ be the center.  Triangulate $H$ by the $30$ triangles
\begin{align*}
&(z,i_k,i_{k+1}),\\
&(i_k,o_k,o_{k+1}),\\
&(i_k,o_{k+1},i_{k+1}),
\end{align*}
with indices modulo $10$.  All squared source lengths lie in $\mathbb Q(\sqrt5)$.

The ten outer vertices are identified in cyclic order as
\begin{equation}\label{eq:boundarywords}
A_2,\ A_1,\ J,\ C,\ J,\ B_1,\ B_2,\ B_1,\ J,\ A_1.
\end{equation}
Thus the boundary traces the graph
\[
A_2-A_1-J-C-J-B_1-B_2-B_1-J-A_1-A_2.
\]
Every seam edge is traversed twice in opposite directions, and the branch $J-C-J$ is the single pleat.  The ten inner vertices and the center remain distinct.  Hence the quotient has $17$ vertices.

\begin{figure}[t]
\centering
\includegraphics[width=.92\textwidth]{figures/construction_multiview.pdf}
\caption{The source triangulation and four opaque views of the certified target surface.  Red edges indicate the seam graph; the points $C$ and $J$ are marked in the target views.  In particular, the opposite, near-side, and top views make clear that $C$ is a vertex of the polyhedral surface, not an interior point.  The drawings are for visualization only; the proof uses the exact coordinates in \cref{sec:coordinates}.}
\label{fig:construction}
\end{figure}

\begin{figure}[H]
\centering
\includegraphics[width=.50\textwidth]{figures/seam_graph.pdf}
\caption{Combinatorial seam graph.  The branch $JC$ is the single pleat.  The actual target coordinates are not coplanar and are listed in \cref{sec:coordinates}.}
\label{fig:seam}
\end{figure}

\section{Explicit target coordinates}\label{sec:coordinates}
Let
\[
\lambda=\frac{99999}{100000}.
\]
Each target vertex is $\lambda$ times the finite-decimal vector in \cref{tab:coords}.  Every displayed decimal is interpreted as an exact rational number.  The affine image of each source triangle is determined by the images of its three vertices.

\begin{table}[H]
\centering
\small
\caption{Base target coordinates before multiplication by $\lambda=99999/100000$.}
\label{tab:coords}
\begin{tabular}{@{}lrrr@{}}
\toprule
vertex & $x$ & $y$ & $z$\\
\midrule
$J$ & 0.114996453431 & 0.652042463578 & 0.359274110539\\
$A_1$ & 0.713527595885 & 0.556758865942 & 0.238247176413\\
$A_2$ & 0.904510611077 & 0.150919970833 & -0.186942534628\\
$B_1$ & -0.269259146678 & 0.427202794161 & 0.787947375499\\
$B_2$ & -0.411232941130 & -0.173302559313 & 0.822628398965\\
$C$ & -0.407447365206 & 0.632905197995 & 0.029649729336\\
$z$ & 0 & 0 & 0\\
$I_0$ & 0.650000000000 & 0 & 0\\
$I_1$ & 0.525861046344 & 0.382060413990 & 0\\
$I_2$ & 0.201227167335 & 0.617682813279 & 0.021807551587\\
$I_3$ & -0.138361314859 & 0.605243587298 & -0.192448295874\\
$I_4$ & -0.192247060016 & 0.586039489637 & 0.205179883277\\
$I_5$ & -0.332154604650 & 0.312020954927 & 0.463482731389\\
$I_6$ & -0.424437564120 & -0.078307526511 & 0.486025395896\\
$I_7$ & -0.126467920004 & 0.156558885543 & 0.618057586773\\
$I_8$ & 0.208351084601 & 0.348929212409 & 0.507285156764\\
$I_9$ & 0.525861046344 & 0.263795278601 & 0.276373318043\\
\bottomrule
\end{tabular}
\end{table}

Let $g:H\to\R^3$ denote the resulting continuous piecewise-affine map.  Continuity holds because adjacent triangles use the same target coordinates on their common vertices, including the identifications in \eqref{eq:boundarywords}.

\section{Exact certification}\label{sec:certificate}
We now establish the geometric properties of $g$.  All sign and intersection assertions in this section are checked by exact arithmetic in the supplied programs; the verification methods are described explicitly so that the certificate is finite and reproducible.

\subsection{Strict nonexpansion on every triangle}
For a source triangle choose two edge vectors $u,v$ from a common vertex and set
\[
S=\begin{pmatrix}
\langle u,u\rangle & \langle u,v\rangle\\
\langle u,v\rangle & \langle v,v\rangle
\end{pmatrix}.
\]
Let $T$ be the corresponding Gram matrix of the target edge vectors.  A tangent vector with coordinate column $\xi\in\R^2$ has squared source length $\xi^TS\xi$ and squared target length $\xi^TT\xi$.  Thus the affine triangle map is strictly nonexpansive if
\[
S-T\succ0.
\]

\begin{lemma}\label[lemma]{lem:lip}
For each of the $30$ source triangles, $S-T$ is positive definite.  Hence $g$ is $1$-Lipschitz on $H$.
\end{lemma}

\begin{proof}
For a real symmetric $2\times2$ matrix, positive definiteness is equivalent to positivity of the two diagonal entries and of the determinant.  The verifier computes, for each triangle,
\[
(S-T)_{11}>0,\qquad (S-T)_{22}>0,\qquad \det(S-T)>0.
\]
The source terms lie in $\mathbb Q(\sqrt5)$ and the target terms are rational; signs in $\mathbb Q(\sqrt5)$ are reduced to rational comparisons by squaring only after the signs of the rational coefficients are known.  Thus no floating-point sign decision is used.

To pass from triangles to the whole decagon, take $x,y\in H$.  Convexity of $H$ implies that the segment $[x,y]$ remains in $H$ and crosses finitely many source triangles.  On each subsegment the image curve has length at most the source subsegment length.  Summing yields
\[
|g(x)-g(y)|\le \operatorname{length}(g([x,y]))\le |x-y|.
\]
\end{proof}

\subsection{The quotient is a sphere}
\begin{lemma}\label[lemma]{lem:sphere}
The quotient triangulation is a connected closed orientable triangulated $2$-manifold with
\[
V=17,\qquad E=45,\qquad F=30,
\]
and hence is topologically $S^2$.
\end{lemma}

\begin{proof}
The exact combinatorial verifier checks that every edge belongs to exactly two faces, every vertex link is one cycle, the complex is connected, and the listed face orientations agree oppositely along each shared edge.  Hence it is a connected closed orientable $2$-manifold.  Its Euler characteristic is
\[
17-45+30=2,
\]
so the classification of closed orientable surfaces gives $S^2$.
\end{proof}

\subsection{No unintended triangle intersections}
\begin{lemma}\label[lemma]{lem:embedded}
The geometric realization of the quotient complex is embedded in $\R^3$.
\end{lemma}

\begin{proof}
There are $\binom{30}{2}=435$ unordered face pairs.  The independent verifier tests intersection by solving exactly
\[
A_0+\alpha(A_1-A_0)+\beta(A_2-A_0)
=
B_0+\gamma(B_1-B_0)+\delta(B_2-B_0)
\]
under the barycentric inequalities
\[
\alpha,\beta,1-\alpha-\beta,\gamma,\delta,1-\gamma-\delta\ge0.
\]
The three coordinate equations have rank $3$ in every tested pair, leaving one affine parameter.  The six inequalities become rational interval constraints on that parameter.

The exhaustive outcome is:
\begin{center}
\begin{tabular}{@{}lrr@{}}
\toprule
prescribed common simplex & number of pairs & exact intersection\\
\midrule
none & 260 & empty\\
one vertex & 130 & that vertex only\\
one edge & 45 & that edge only\\
\bottomrule
\end{tabular}
\end{center}
Thus there are no unintended self-intersections.  Together with \cref{lem:sphere}, the realization is an embedded PL $2$-sphere.  By Jordan--Brouwer separation it bounds a compact solid, denoted $B_1$.
\end{proof}

\subsection{Nonconvexity}
\begin{lemma}\label[lemma]{lem:nonconvex}
The bounded solid $B_1$ is nonconvex.
\end{lemma}

\begin{proof}
Consider the face $(z,I_0,I_1)$ and let $n$ be its oriented normal.  Exact rational arithmetic gives
\[
n\cdot(J-z)=
\frac{178438386772277185863216936460667247207}
{2000000000000000000000000000000000000000}>0,
\]
whereas
\[
n\cdot(A_2-z)=
-\frac{23211889835379489551564573316040740141}
{500000000000000000000000000000000000000}<0.
\]
Thus vertices of the boundary lie on both sides of the plane of a boundary triangle.  Every face plane of a convex polyhedron is supporting, so this is impossible for a convex solid.
\end{proof}

\subsection{Exact volume}
\begin{lemma}\label[lemma]{lem:volume}
The volume enclosed by the target sphere equals the constant $C$ in \eqref{eq:Cexact}.
\end{lemma}

\begin{proof}
The face orientations are globally consistent.  Therefore the oriented tetrahedral formula gives
\[
\Vol(B_1)=\frac16\left|\sum_{(a,b,c)\in\mathcal F}\det(a,b,c)\right|.
\]
All target coordinates are rational.  Exact evaluation yields
\[
\Vol(B_1)=
\frac{
359686988325832312742347877865122384142690317017149
}{
2000000000000000000000000000000000000000000000000000
},
\]
which is $0.17984349416291615\ldots>0.179843$.
\end{proof}

\section{From the decagon to the disk}\label{sec:disk}
The preceding construction uses only the regular decagon $H$, but the statement concerns the entire disk.  The transfer is immediate from metric projection.

\begin{lemma}\label[lemma]{lem:projection}
Let $\pi_H:D_1\to H$ be Euclidean nearest-point projection.  Then $\pi_H$ is $1$-Lipschitz and onto.  Consequently
\[
F_1:=g\circ\pi_H:D_1\to\R^3
\]
is $1$-Lipschitz and has image $\partial B_1$.
\end{lemma}

\begin{proof}
For $x,y\in D_1$, set $p=\pi_H(x)$ and $q=\pi_H(y)$.  The standard projection inequalities for a closed convex set are
\[
(x-p)\cdot(q-p)\le0,
\qquad
(y-q)\cdot(p-q)\le0.
\]
Adding gives
\[
|p-q|^2\le(p-q)\cdot(x-y)\le |p-q|\,|x-y|,
\]
and hence $|p-q|\le|x-y|$.  The projection fixes every point of $H$, so it is onto.  Therefore $F_1$ is $1$-Lipschitz and
\[
F_1(D_1)=g(H)=\partial B_1.
\]
\end{proof}

\begin{lemma}[Boundary self-gluing]\label[lemma]{lem:boundarygluing}
The restriction $\pi_H|_{\partial D_1}:\partial D_1\to\partial H$ is onto.  Under $F_1=g\circ\pi_H$, every relative interior point of a seam edge has at least two distinct preimages on $\partial D_1$, coming from the two oppositely oriented boundary edges in \eqref{eq:boundarywords}.  Thus no boundary arc remains unglued.
\end{lemma}

\begin{proof}
Let $p$ lie in the relative interior of an edge $e$ of the inscribed decagon $H$, and let $n_e$ be its outward unit normal.  Because $p$ lies strictly inside the unit circle, the ray $p+t n_e$, $t\ge0$, meets $\partial D_1$ at some point $x$.  Since $x-p$ lies in the normal cone of the convex set $H$ at $p$, nearest-point projection satisfies $\pi_H(x)=p$.  At a vertex $p=o_k$ we simply have $p\in\partial D_1$ and $\pi_H(p)=p$.  Hence $\pi_H(\partial D_1)=\partial H$.

The cyclic word \eqref{eq:boundarywords} traverses each of the five seam segments $A_2A_1$, $A_1J$, $JC$, $JB_1$, and $B_1B_2$ exactly twice, once in each direction.  On each outer decagon edge, $g$ is affine and maps that edge onto the corresponding seam segment.  Therefore a relative interior point $y$ of any seam segment has preimages $p,p'$ in the relative interiors of two distinct decagon edges.  By the first paragraph there are $x,x'\in\partial D_1$ with $\pi_H(x)=p$ and $\pi_H(x')=p'$, so $x\ne x'$ and $F_1(x)=F_1(x')=y$.  At the terminal seam vertices $A_2$, $C$, and $B_2$, the two paired boundary arcs meet at the same image vertex; at the intermediate vertices $A_1$, $J$, and $B_1$, the incident paired arcs meet according to the cyclic word \eqref{eq:boundarywords}.  In every case the corresponding quotient vertex link, checked in \cref{lem:sphere}, is a circle rather than an interval.  Thus the relative interiors are paired and the endpoints close up without leaving boundary, so the whole disk boundary is self-glued in the image.
\end{proof}

Scaling both the planar domain and the target by $R$ defines $F_R$ and gives $\Vol(B_R)=R^3\Vol(B_1)$.  Together, \cref{lem:lip,lem:sphere,lem:embedded,lem:nonconvex,lem:volume,lem:projection,lem:boundarygluing} prove \cref{thm:main}.

\section{Verification protocol and reproducibility}\label{sec:repro}
The accompanying repository contains two independent exact verifiers together with a \texttt{README.md} giving the exact commands, dependencies, and expected certificate output.

\subsubsection*{Verifier A}
\texttt{certificate/exact\_certificate.py}.\par
Target coordinates use Python's \texttt{fractions.Fraction}; source metric expressions use exact symbolic arithmetic in $\mathbb Q(\sqrt5)$.  Triangle intersections are tested by exact plane-section interval arithmetic.

\subsubsection*{Verifier B}
\texttt{certificate/independent\_verify.py}.\par
This implementation does not use SymPy.  It contains a hand-written $\mathbb Q(\sqrt5)$ arithmetic type, checks orientability by propagation on the dual graph, and tests triangle intersections by the barycentric feasibility method described in \cref{lem:embedded}.

Both programs independently return the same exact rational volume and the same intersection counts
\[
260+130+45=435.
\]
A successful run prints, among other diagnostics,
\begin{verbatim}
V,E,F = 17 45 30 Euler = 2
triangle pairs = {1: 130, 2: 45, 0: 260}
volume decimal = 0.17984349416291615
all 30 triangle maps are strictly 1-Lipschitz
\end{verbatim}

The scale factor $\lambda=0.99999$ is intentionally conservative.  It creates a strict nonexpansion margin while preserving the combinatorial and intersection structure under uniform scaling.  The theorem makes no optimality claim about the displayed solid.

\section{Discussion}
The main conclusion is the certified lower bound
\[
W(R)\ge C R^3,
\qquad
C=0.17984349416291615\ldots .
\]
The construction does not identify $W(R)$ and makes no claim that one pleat is optimal.  It also does not determine the supremum restricted to convex wrappable bodies.  The Mylar comparison is therefore deliberately modest: it shows that the certified witness is quantitatively well beyond the sharp convex body-of-revolution benchmark, not that it beats every convex body wrappable by a disk.

The soft argument in \cref{sec:soft} changes the interpretation of Nandakumar's Question~1.  Under the stated $1$-Lipschitz formalization, nonconvexity by itself is easy to introduce once one has a rotational wrap with strict meridian-length and volume slack.  A tiny inward notch already gives an affirmative answer.  The more substantive extremal question is how much volume can be enclosed, and the exact coefficient $C$ should be viewed as a concrete lower bound for that problem rather than as a certificate needed merely to establish nonconvex existence.

The use of a piecewise-affine source map is distinct from the problem of finding a neat isometric origami folding.  Work of Demaine and Ku shows that, for polygonal boundaries folded nonexpansively at finitely many points, compatible isometric mappings of the polygon interior exist under their model~\cite{DemaineKu2014}.  The present proof does not invoke that theorem: the full non-neat $1$-Lipschitz map is written down directly and independently certified for self-intersection.

Natural next problems are to optimize the present crease network, to search for larger certified coefficients using finer source complexes, and to understand the best convex body wrappable by a disk.  Any improvement of the first two types would raise the explicit lower bound on $W(1)$; the convex problem would clarify how much of the gain is genuinely tied to nonconvexity.

\section*{Acknowledgements}
The coordinate configuration was found by computer-assisted geometric search.  The final proof was reorganized into a finite exact certificate and checked by two independent implementations.  External verification of both the mathematical interpretation and the code is encouraged.

\begin{thebibliography}{9}
\bibitem{Nandakumar2026}
R.~Nandakumar,
\newblock \emph{On the maximum volume solid wrappable by a given sheet of paper},
\newblock arXiv:2604.02925v1 [math.MG], 2026.

\bibitem{NandakumarMO2026}
R.~Nandakumar,
\newblock \emph{A claim on wrapping 3D solids with a sheet of paper},
\newblock MathOverflow question 508997, 2026; remark added April 21, 2026.

\bibitem{Paulsen1994}
W.~H. Paulsen,
\newblock \emph{What Is the Shape of a Mylar Balloon?},
\newblock Amer. Math. Monthly \textbf{101} (1994), no.~10, 953--958.

\bibitem{MladenovOprea2003}
I.~M. Mladenov and J.~Oprea,
\newblock \emph{The Mylar Balloon Revisited},
\newblock Amer. Math. Monthly \textbf{110} (2003), no.~9, 761--784.

\bibitem{DemaineKu2014}
E.~D. Demaine and J.~S. Ku,
\newblock \emph{Filling a Hole in a Crease Pattern: Isometric Mapping of a Polygon given a Folding of its Boundary},
\newblock in \emph{Origami$^6$: Proceedings of the 6th International Meeting on Origami in Science, Mathematics and Education}, American Mathematical Society, pp.~177--188, 2014; arXiv:1410.6520.
\end{thebibliography}

\end{document}
